\item En un experimento biológico de absorción de nutrientes en plantas se emplean dos radioisótopos, $X$ e $Y$. Inicialmente hay $2.50$ veces más núcleos de tipo $X$ que de tipo $Y$. Transcurridos $3.00\text{ días}$, la proporción cambia de modo que hay $4.20$ veces más núcleos de tipo $X$ que de tipo $Y$. Sabiendo que el isótopo $Y$ tiene una vida media de $1.60\text{ días}$, calcule la vida media del isótopo $X$.